% Generated draft: requires mathematical and source review.
\documentclass[11pt]{article}
\usepackage[margin=25mm]{geometry}
\usepackage{fontspec}
\setmainfont{Latin Modern Roman}
\usepackage{amsmath,amssymb,mathtools,unicode-math}
\setmathfont{Latin Modern Math}
\usepackage{enumitem,needspace,xurl,hyperref}
\setlength{\parindent}{0pt}
\setlength{\parskip}{5pt}
\setlength{\emergencystretch}{3em}
\newcommand{\N}{\mathbb{N}}
\newcommand{\Z}{\mathbb{Z}}
\newcommand{\Q}{\mathbb{Q}}
\newcommand{\R}{\mathbb{R}}
\newcommand{\C}{\mathbb{C}}
\newcommand{\F}{\mathbb{F}}
\newcommand{\A}{\mathbb{A}}
\newcommand{\PP}{\mathbb{P}}
\newcommand{\E}{\mathbb{E}}
\newcommand{\eps}{\varepsilon}
\newcommand{\dd}{\,\mathrm{d}}
\begin{document}
\section*{UnsolvedMath ID: 30005424}
\textbf{Source problem code:} catalogue entry
\subsection{Definitions and mathematical statement}
For a ring \(R\), use the convention that \(\operatorname{mod}R\) denotes the category of finitely generated \(R\)-modules. Let \(Q\) be the quiver with vertices \(1,2\) and a single arrow \(1\to2\). Define \(\operatorname{mono}(Q,\operatorname{mod}R)\) to have objects the injective \(R\)-linear maps \(u\colon M_1\hookrightarrow M_2\). A morphism from \(u\colon M_1\hookrightarrow M_2\) to \(v\colon N_1\hookrightarrow N_2\) is a pair of \(R\)-linear maps \((a,b)\), where \(a\colon M_1\to N_1\), \(b\colon M_2\to N_2\), and \(b\circ u=v\circ a\); composition is componentwise.

With the conventional interpretation that \(p\) is a prime and \(k\) is a field, put
\[
\mathcal C_p=\operatorname{mono}\bigl(Q,\operatorname{mod}(\Z/p^6\Z)\bigr),\qquad
\mathcal D_k=\operatorname{mono}\bigl(Q,\operatorname{mod}(k[x]/(x^6))\bigr).
\]
The supplied excerpt does not explicitly specify these parameter domains, the permitted fields \(k\), any relation between \(p\) and \(k\), or whether the comparison is intended for all parameters or for particular choices. The finite-generation convention above also requires confirmation against the complete source.

\textbf{Comparison problem.} Compare the representation theory of \(\mathcal C_p\) and \(\mathcal D_k\). In the wording of the supplied original excerpt, is their structure similar? The cleaned registry statement asks whether they are equivalent in an appropriate representation-theoretic sense. Neither wording specifies a precise relation to be established.

The excerpt separately defines, attributing the definition to Auslander, a \emph{representation equivalence} to be a full, dense functor \(F\) that reflects isomorphisms. Here dense means that every target object is isomorphic to an object in the image of \(F\), and reflection of isomorphisms means
\[
F(f)\text{ is an isomorphism}\ \Longrightarrow\ f\text{ is an isomorphism}.
\]
This supplies a possible comparison notion, but the excerpt does not identify it as the required relation between \(\mathcal C_p\) and \(\mathcal D_k\), specify a direction for such a functor, or state whether ordinary categories or injectively stable categories are intended.

The surrounding discussion assumes an abelian category \(\mathcal A\) with enough injectives. It describes \(\operatorname{mono}(Q,\mathcal A)\) as an extension-closed subcategory of \(\operatorname{rep}(Q,\mathcal A)\), equipped with its induced exact structure and enough injective objects. It distinguishes these exact-category injectives from injectives in the ambient representation category. An injectively stable category uses the quotient by morphisms factoring through the relevant injective objects, so that distinction matters if a stable comparison is intended. The supplied excerpt loses notation distinguishing stable categories and ends mid-sentence; its contextual theorem cannot therefore resolve the intended comparison target reliably.

Consequently, this is a statement-only draft of the supplied comparison problem. A uniquely quantified equivalence assertion requires clarification of the parameters, comparison relation, and categories to which it applies.
\subsection{Short English statement}
Compare categories of embeddings of finitely generated modules over \(\Z/p^6\Z\) and \(k[x]/(x^6)\). The source asks whether their representation-theoretic structures are similar. It does not specify the relationship between \(p\) and \(k\), the parameter quantifiers, or the precise meaning of similarity or equivalence. The supplied open status is dated curation, not an independently verified conclusion.
\subsection{Sources}
\begin{itemize}
\item Supplied local UnsolvedMath registry record, permanent ID \(30005424\), title \emph{Comparison of Integral and Truncated-Polynomial Birkhoff Categories}, source problem code \texttt{catalogue entry}, in the Oberwolfach Reports Open Problems collection.
\item \emph{Representation Theory of Quivers and Finite-Dimensional Algebras} (2023), Oberwolfach Reports, pp.\ 397--486, DOI \texttt{10.4171/owr/2023/7}. Supplied report link and literature-source entry labelled ``original OWR report'': \url{https://doi.org/10.4171/owr/2023/7}. This report citation is supplied authoritative provenance; it has not been independently consulted for this draft.
\item Supplied literature-source entry: Gao--K\"ulshammer--Kvamme--Psaroudakis, PLMS (2024), \url{https://londmathsoc.onlinelibrary.wiley.com/doi/full/10.1112/plms.12640}. The record supplies no article title or further bibliographic particulars.
\item Attribution in the supplied excerpt: a question is attributed to Birkhoff's paper labelled [1]; the structural comparison is attributed to the authors of a work labelled [7]; the definition of representation equivalence is attributed to Auslander. The bibliography entries corresponding to [1] and [7] are not supplied, and the referent of the initial phrase about Birkhoff is incomplete. No proposer is named in the registry.
\item Literature assessment checked at \texttt{2026-08-22}: the supplied curation labels the problem open and describes it as a boundary-case comparison, says the truncated-polynomial category is much better understood, and reports that no equivalence or matching structure has been proved. These are supplied dated curation claims, not independently verified literature conclusions. The background explicitly requires independent verification of time-sensitive status and cited claims.
\item The registry labels the statement \texttt{corrected\_verified} and reports a check against the official Oberwolfach report. This is a supplied verification claim, not a verification performed for this draft.
\end{itemize}
\end{document}
